% =============================================================================
\chapter{Knowledge suppliers as error processes}\label{ch:suppliers}

Chapter~\ref{ch:frame} built the rows and the questions; this chapter is about who answered them.
The question is not whether a language model understands causality. It is whether the wrong claims the
certificate was tested against came from processes varied and realistic enough that the verdict of
Chapter~\ref{ch:table4} does not depend on one way of being wrong. The models are suppliers of claims,
one of four knowledge processes, and nothing in the paper ranks them.

\begin{ideia}
A certificate against wrong claims can only be tested with wrong claims, and whoever chooses
\emph{which} claims are wrong decides what the test can see. So the evaluation uses several error
processes, including one whose errors nobody chose.
\end{ideia}

% -----------------------------------------------------------------------------
\begin{antes}
\antesq[fornecedor-conhecimento]{Which four processes produce the wrong claims of the evaluation, and
what does each answer that the others cannot?}
\antesq[fornecedor-conhecimento]{Why are commission and omission two columns and never a sum?}
\end{antes}

\section{Four sources of wrong claims}\label{sec:sup-sources}

A synthetic reversal chooses the error: the experimenter knows how many claims are wrong and where,
which is what the knowledge interval needs, so the controlled arms $b = 0, 1, 2$ feed Table~1
(\texttt{results/interval/PREREGISTRATION.md}; Chapter~\ref{ch:gamma}). But a reversal never decides
what to talk about, never correlates neighbouring errors and never declines. A real source does all
three. The evaluation therefore has four processes, each answering one question:
\begin{itemize}[nosep, leftmargin=*]
\item controlled reversals: what exactly $b$ wrong claims cost;
\item chance (\texttt{D\_RAND}): orientations drawn among the Meek-consistent choices on the
  primary supplier's own edges, that is, what a coin does where the source spoke;
\item language-model elicitation: an uncontrolled error process that chooses what to assert,
  what to decline and where to err;
\item truthful controls (\texttt{A\_TRUE}, and \texttt{A\_ORACLE}), read off the truth and
  printed below the rule: where one asserts, separated from asserting wrongly.
\end{itemize}
An arm with no knowledge (\texttt{D\_DEGEN}) is the floor.

\begin{sloppypar}
The local suppliers ran at temperature zero and seed zero, each answer cached under the SHA-256 of
model and prompt (\texttt{experiments/elicit\_run.py}). The primary supplier,
\texttt{qwen2.5:7b-instruct}, answered on real names (\texttt{D\_LLM}), on the second presentation order
(\texttt{D\_LLM\_INSTR}) and on scrambled names (\texttt{D\_SCRAMBLED}); a second family,
\texttt{llama3.1:8b}, on real names (\texttt{D\_LLM\_SRC}). A twenty-billion reasoning model on the same
host spent its whole budget reasoning, returned nothing, and was dropped.
\end{sloppypar}

\par\ifdim\dimexpr\pagegoal-\pagetotal\relax<4.2cm\newpage\fi
\begin{exemplo}[What the primary supplier did with 369 open pairs]
Every open pair of the questionnaire ends in exactly one box.
\begin{center}
\begin{tikzpicture}[x=1cm, y=1cm, font=\scriptsize\sffamily]
\node[nodot] (all) at (0,-0.3) {369 asked};
\node[nodot, densely dashed] (dec) at (2.3,0.3) {137 declined};
\node[nodot, densely dashed] (nr)  at (2.3,-0.3) {18 not reached};
\node[nodot] (ass) at (2.3,-0.9) {214 claims made};
\node[nodot, densely dashed] (drop) at (5.4,-0.6) {43 dropped by the repair};
\node[nodot] (kept) at (5.4,-1.2) {171 kept};
\node[nodot, draw=verdesc, text=verdesc] (tru) at (8.5,-0.9) {118 true};
\node[nodot, draw=vermesc, text=vermesc] (fal) at (8.5,-1.5) {\textbf{53 false} (commission)};
\foreach \a/\b in {all/dec, all/nr, all/ass, ass/drop, ass/kept, kept/tru, kept/fal}
  \draw[cedge] (\a.east) -- (\b.west);
\end{tikzpicture}
\end{center}
Dashed boxes end with no kept claim: $137 + 18 + 43 = 198$ pairs, the omission. Counts from
\texttt{TABLE4.md} (elicitation block) and \texttt{knowledge\_audit.json}.
\end{exemplo}

The tree separates two quantities, and the protocol names them so that nobody adds them.

\begin{definicao}[Commission and omission, at the knowledge level]
For one supplier and one network, \emph{commission} is the fraction of claims kept after the repair
that are false at the truth, and \emph{omission} is the fraction of asked pairs that end with no kept
claim: declined, not reached or dropped. Both are averaged over networks with a network bootstrap and
printed in separate columns (\texttt{experiments/knowledge\_audit.py}).
\end{definicao}

A decline leaves the edge open and the analyst's graph wider; a false claim closes the edge the wrong
way and can invalidate the set. The denominators also differ, so the sum is not the fraction of
anything.

\begin{medido}[Primary supplier, \texttt{results/elicit/knowledge\_audit.json}]
Commission \textbf{0.325} [0.210, 0.449] by network mean (53 of 171 pooled); omission 0.450 [0.361,
0.545]. By tier: 0.27 on the benchmarks, 0.28 on the applied-paper diagrams, 0.44 on the four Gaussian
networks, \textbf{0.75} on the textbook diagrams, whose variables are abstract symbols. L1 registered a band of
[0.10, 0.45]: supported.
\end{medido}

The protocol's two other suppliers, claims harvested from the source publications and human analysts,
are absent: both cost human hours, and \texttt{docs/PROTOCOL\_RUN.md} records the deferral with that
reason.

\prompt[conceito=fornecedor-conhecimento,dificuldade=0.7]{The primary supplier has commission 0.325
and omission 0.450. Why is there no ``total error'' column?}{The two act in opposite directions on the
certificate and have different denominators. A decline leaves an edge open, widens the analyst's graph
and may make the query non-amenable; a false claim closes the edge the wrong way and may invalidate
the committed set. Commission is a fraction of kept claims, omission a fraction of asked pairs, so their
sum is a fraction of nothing.}

% -----------------------------------------------------------------------------
\begin{antes}
\antesq[escala-validade]{Of the predictions S1 to S10, which failure changes what the paper says?}
\antesq[escala-validade]{How can a supplier that is wrong less often per claim have its committed set
invalid as often?}
\end{antes}

\section{Scale buys claims, not validity}\label{sec:sup-scale}

On 12 September four checkpoints on one node of four H100 answered the same 540 frozen items through
vLLM, with each model's chat template, temperature zero, seed zero and 700 tokens:
\texttt{Qwen2.5-72B-Instruct} (real names at both orders, and scrambled), \texttt{Qwen2.5-32B-Instruct},
\texttt{Llama-3.3-70B-Instruct} and \texttt{gemma-3-27b-it}. The 2,160 answers entered the same cache
and went through the same parser, repair and sweep: 8,308 ledger rows, 3,014 certifiable
rows on the true CPDAG. Predictions S1 to S10 were written after the jobs were submitted and before any answer was read,
with a decision rule: if S1 and S6 both held, the 72B row would lead.

\begin{center}
\footnotesize
\begin{tabular}{@{}l>{\raggedright\arraybackslash}p{4.8cm}>{\raggedright\arraybackslash}p{3.15cm}>{\raggedright\arraybackslash}p{6.4cm}@{}}
\toprule
& prediction & verdict & number \\
\midrule
S1 & scale lowers commission & not as registered & 72B $-$ 7B: $-0.100$ [$-0.224$, $+0.017$] \\
S2 & scale lowers omission & supported & 72B $-$ 7B: $-0.280$ [$-0.390$, $-0.164$] \\
S3 & at 72B, names carry information & not supported & scrambled $-$ real $+0.088$ [$-0.064$, $+0.238$] \\
S4 & claim certificate never dangerous & supported, structural & 0 of 3,014 certifiable rows \\
S5 & the unit gap does not close & supported, thinly & hop: 1 dangerous row at 72B, 0 on gemma \\
S6 & scale lowers invalid sets & \textbf{not supported} & 0.298 at 72B, 0.288 at 7B \\
S7 & moves do not order by truth & unresolved, as predicted & 0.846 at 72B, 1.086 on the truthful arm \\
S8 & the third family is no outlier & not supported & gemma $-$ 72B: $+0.112$ [$+0.029$, $+0.213$] \\
S9 & less source variance at scale & direction only & 0.182 against 0.220 \\
S10 & bridge improves, stays biased & not supported & 72B: control $-$ open $+0.004$ [$-0.153$, $+0.159$] \\
\bottomrule
\end{tabular}
\end{center}

S6 is the row that decides (\texttt{results/elicit/PREREGISTRATION\_SCALE.md}, Appendix~A). The 72B
keeps claims on 281 of the 369 pairs and declines 27, where the 7B kept 171 and declined 137. Commission
falls from 0.325 to 0.233 by network mean without a separated paired interval, while the committed set
is invalid at the truth on 85 of 285 certifiable rows at 72B and on 66 of 229 at 7B. S1 did not separate
and S6 failed, so the 7B stays the primary supplier and the 72B is the scale replicate.
\figref{fig:supplier-scale} puts all ten conditions side by side.

\begin{figure}[htbp]
\centering
\includegraphics[width=\linewidth,height=0.25\textheight,keepaspectratio]{supplier_scale.pdf}
\caption{Scale buys claims and accuracy per claim, not validity: across the Qwen family the number of
kept claims rises from 171 to 223 to 281 and the false fraction falls, while the share of certifiable
rows whose committed set is invalid at the truth goes 0.288, 0.201, 0.298, in no order and with
overlapping intervals.}
\label{fig:supplier-scale}
\end{figure}

The mechanism is counting. A set is invalid when one wrong claim sits where the query needs it, and a
supplier that asserts more puts more claims near each query. On commission, gemma is the high outlier
against the 72B.
At 72B the model is also far more stable under presentation order: instrument disagreement
\textbf{0.047} [0.019, 0.079] on shared edges, against 0.160 at 7B (\texttt{TABLE4.md}). The batch
script is \texttt{experiments/remote/elicit\_jz.slurm} (one node, four H100, every cache on scratch
because vLLM's compile cache overflowed the home quota, eager mode); \texttt{elicit\_import.py} puts the
answers under the local cache key, and \texttt{elicit\_run.py} refuses to call a cluster model.


\prompt[conceito=escala-validade,dificuldade=0.7]{Without numbers: why can a supplier that is wrong
less often per claim leave the committed set invalid just as often?}{A committed set is invalid when at
least one wrong claim sits where the query needs it. The larger model declines far less, so it asserts
more edges in each component and more queries become certifiable; every extra claim near a query is
another chance that one of them is wrong. A lower rate per claim times more claims per query leaves the
chance of at least one wrong claim where it was. Accuracy per claim and validity per query are
different quantities, and the certificate sits between them.}

% -----------------------------------------------------------------------------
\begin{antes}
\antesq[sobreposicao-fonte]{Asked cold for the publication behind each network, how many do the local
suppliers name?}
\antesq[ponte-compelida]{Is the error on compelled edges, which needs no truth, an unbiased estimate of
commission on open edges?}
\end{antes}

\section{Names, a truth-free bridge, and two negative results}\label{sec:sup-names}

\emph{Scrambled names.} At 7B, scrambling gives commission 0.343 [0.260, 0.431] against 0.325; what
changes is declining, omission 0.568 against 0.450 (L3 not supported). At 72B the gap is $+0.088$
[$-0.064$, $+0.238$] and the model declines 96 pairs against 27 (S3). A source reciting the benchmarks
would collapse without the names, and neither does; by invariance, that licenses ``no recitation
detected'' and nothing stronger.

\emph{The compelled-edge bridge.} The control block asks about edges the data already orient, so its
error needs no truth. If it estimated open-edge commission without bias, a tier with no known truth
could be scored. At 7B it does not: 0.191 [0.094, 0.293] against 0.325, a correction of about 0.13 that
varies by network (L9 not supported).

\begin{medido}[Control minus open-edge commission, \texttt{results/elicit/knowledge\_audit.json}]
Paired by network: 7B \textbf{$-0.153$} [$-0.324$, $+0.009$] on 23 networks; 72B \textbf{$+0.004$}
[$-0.153$, $+0.159$] on 26. The 72B's own control error is 0.232 [0.121, 0.358], higher than the 7B's
0.191, so size did not improve the control; it removed the bias. At 7B the paired interval reaches
$+0.009$: the bias rests on the registered comparison and a point estimate of $-0.15$.
\end{medido}

The bridge is therefore reported per supplier: usable at 72B, biased at 7B.

\emph{The source-overlap probe.} The scrambled arm answers the recitation question only by invariance,
so a second probe asks it directly: shown the variable list alone, name the publication behind the
network and its declared exposure, outcome and adjustment set. ``Recited'' requires the source and the
query.

\begin{medido}[Source overlap, \texttt{results/elicit/SOURCE\_OVERLAP.md}]
32 networks, 29 with a scorable source, 10 with a declared query. Source named: \textbf{0 of 29} for
\texttt{qwen2.5:7b}, which declines on 28 and misattributes the one it answers; \textbf{0 of 29} for
\texttt{llama3.1:8b}, which answers all 29 and matches none. Declared query reproduced on 2 of 10 by
each, never with the source. \textbf{Recited: none}, for both.
\end{medido}

Matching a query without the source is reading the names. The pre-registration's appendix scores an
earlier pass on prompts rendered with the salted seed of Chapter~\ref{ch:frame} (27 declines, same
verdict); the box is the rerun on stable prompts. The cluster models have not been probed.

\emph{Ancestral claims.} Asked about 156 non-adjacent pairs at skeleton distance two, the 7B asserted
ancestry on \textbf{6}, said ``neither'' on 45 and declined 105; the six reduce soundly to ten
orientations, four false, which change one true claim on one swept network and no certificate
(\texttt{results/elicit/PREREGISTRATION\_ANCESTRAL.md}). All three predictions fell: a negative result.

\prompt[conceito=ponte-compelida,dificuldade=0.75]{At which supplier size is the compelled-edge bridge
unbiased, what is the number, and why does that not mean the larger model is better on compelled
edges?}{At 72B: control minus open-edge commission $+0.004$ [$-0.153$, $+0.159$], paired over 26
networks, against $-0.153$ [$-0.324$, $+0.009$] at 7B. The 72B's error on compelled edges is 0.232,
higher than the 7B's 0.191; the bridge is unbiased because its control error and its open-edge
commission now agree, not because the control error fell.}

% -----------------------------------------------------------------------------
\begin{antes}
\antesq[teste-cascata]{Does the greedy recovering set, which the old pipeline used as knowledge,
orient more edges per claim than the 7B supplier?}
\antesq[teste-cascata]{Which justification for the redesign was withdrawn, and what still carries it?}
\end{antes}

\section{The cascade test: a justification withdrawn}\label{sec:sup-cascade}

The committed pipeline's knowledge was the greedy recovering set, now \texttt{A\_ORACLE}: few claims,
chosen from the truth so that Meek's closure recovers as much as possible. The redesign had two reasons
to drop it: the invariant, and ``the old generator cascades unrealistically'', registered to be
withdrawn unless the oracle set compels more edges per claim than an elicited set, paired by network.

\begin{exemplo}[Three sources on \texttt{child}]
The 7B keeps one claim, the false \texttt{CO2} $\to$ \texttt{LungParench} of Chapter~\ref{ch:frame},
and the closure orients 10 edges the CPDAG left open. The oracle set asserts 2 and orients 12; the 72B
asserts 8 and orients 9. Per claim: 10, 6 and 1.125. The same edge oriented truthfully,
\texttt{LungParench} $\to$ \texttt{CO2}, orients 2 (\texttt{results/ledger/rows.csv}, columns
\texttt{n\_k} and \texttt{k\_g0}). On this network the 7B cascades most.
\end{exemplo}

The quantity is \emph{closure per claim}, $k_{\Gz}/|\Kset|$: edges the analyst's graph orients that
$\Chat$ left open, asserted or propagated, over kept claims. It is constant within a network and an arm,
so the test pairs arms by network (\texttt{experiments/cascade\_test.py}).

\begin{medido}[Cascade test, \texttt{results/ledger/cascade\_test.json}]
Oracle $-$ 7B: \textbf{$+1.13$} [$-0.26$, $+2.86$] on 27 networks, larger on 14: \textbf{not
separated}. Oracle $-$ 72B: $+1.28$ [$+0.43$, $+2.50$]; oracle $-$ Llama 70B: $+1.62$ [$+0.32$,
$+3.27$]; both separated. Truthful orientations on the 7B's edges $-$ 7B: $-0.63$ [$-1.52$, $-0.02$].
Network means: oracle 2.97, 7B 1.84, 70B 1.35, 72B 1.15.
\end{medido}

The justification is withdrawn as written (\texttt{results/ledger/CHECKS.md}). \figref{fig:cascade}
shows why ``unrealistic'' had no fixed referent: the 7B declines a lot, keeps few claims, and each drags
others along, as the oracle's minimal set does; the large models assert almost every edge of a
component, and then the closure has little left to add.

\begin{figure}[htbp]
\centering
\includegraphics[width=\linewidth,height=0.22\textheight,keepaspectratio]{cascade.pdf}
\caption{Whether the old generator ``cascades too much'' depends on the supplier it is compared with:
the oracle set does not separate from the 7B supplier in edges oriented per claim, and does separate
from the 72B and 70B suppliers, which assert many claims that propagate little. The redesign keeps
only the invariant as its reason.}
\label{fig:cascade}
\end{figure}

\prompt[conceito=teste-cascata,dificuldade=0.7]{Explain without contradiction why the same oracle set
does not cascade more than the 7B supplier but does cascade more than the 72B.}{Closure per claim
depends on how much a source says. The 7B declines a lot (137 of 369 pairs) and keeps few claims per
component, and each drags orientations along through Meek's rules, as the oracle's minimal set does.
The 72B keeps claims on almost every pair (281 of 369); when nearly every edge of a component is asserted
directly, the closure adds little and the ratio is close to 1. ``Unrealistic'' only means something
against a reference source, and the real sources disagree with each other.}

\naopasse[escala-validade]{Write the chain from ``the 72B keeps 281 claims'' to ``its set is invalid
on 0.298 of certifiable rows''. If the step from more claims to more invalid sets does not come, reread
Section~\ref{sec:sup-scale}.}

\begin{exercicio}[metodo=table-reading,conceito=escala-validade]
From \texttt{TABLE4.md} and \texttt{knowledge\_audit.json}, the six real-name suppliers at the first
order. (a) Within the Qwen
family, does the invalid-set rate fall with size? (b) Which column explains gemma? (c) Check that the
first four columns account for all 369 pairs for the 72B and for Llama 70B. (d) Which supplier is
``best'', and why does that question not belong in the paper?

\begin{center}
\small
\begin{tabular}{@{}lrrrrrrr@{}}
\toprule
supplier & kept & dropped & declined & not reached & commission & certifiable & invalid \\
\midrule
Qwen 7B   & 171 & 43 & 137 & 18 & 0.325 & 229 & 0.288 \\
Qwen 32B  & 223 & 26 &  83 & 37 & 0.271 & 204 & 0.201 \\
Qwen 72B  & 281 & 24 &  27 & 37 & 0.233 & 285 & 0.298 \\
Llama 8B  & 223 & 41 &  92 & 13 & 0.308 & 243 & 0.230 \\
Llama 70B & 282 & 78 &   9 &  0 & 0.259 & 267 & 0.247 \\
gemma 27B & 194 & 68 &  14 & 93 & 0.324 & 263 & 0.308 \\
\bottomrule
\end{tabular}
\end{center}
\end{exercicio}
\begin{solucao}
(a) No: 0.288, 0.201, 0.298, in no order, and Table~4's intervals overlap. (b) ``Not reached'': gemma
leaves 93 pairs unanswered, drops 68 in the repair and keeps a commission as high as the 7B's.
(c) 72B: $281 + 24 + 27 + 37 = 369$. Llama 70B: $282 + 78 + 9 + 0 = 369$. (d) None: the suppliers are
error sources for testing instruments, and what the paper uses is that validity does not improve with
scale.
\end{solucao}

\begin{exercicio}[metodo=estimation,conceito=escala-validade]
Model a certifiable row as depending on $m$ claims, each false independently with probability equal to
the supplier's commission $c$, so the set is invalid with probability $1 - (1-c)^m$. (a) Estimate the
effective $m$ that reproduces the invalid-set rate for the 7B ($c = 0.325$, rate 0.288) and for the 72B
($c = 0.233$, rate 0.298). Use $\ln 0.675 \approx -0.393$, $\ln 0.712 \approx -0.340$,
$\ln 0.767 \approx -0.265$, $\ln 0.702 \approx -0.354$. (b) Compare the ratio of the two $m$ with the
ratio of pairs the two suppliers kept. (c) Name two reasons the model is too crude to report.
\end{exercicio}
\begin{solucao}
(a) 7B: $m = \ln 0.712 / \ln 0.675 \approx 0.340/0.393 \approx 0.86$. 72B: $m = \ln 0.702 / \ln 0.767
\approx 0.354/0.265 \approx 1.33$. (b) $1.33/0.86 \approx 1.5$, against $281/171 \approx 1.6$ kept
claims: the rise in effective claims per query is about the rise in claims, which is the counting
mechanism in numbers. (c) Errors are not independent (one answer, one repair, clustering by network); validity
depends on the claims on the query's paths, not on a random $m$; and the rates come from different
certifiable rows, 229 and 285. An order-of-magnitude check, not a result.
\end{solucao}

\section*{Chapter map}

\begin{center}
\begin{tikzpicture}[x=1cm, y=1cm]
\node[cmap]  (q)  at (0, 0)      {four knowledge processes};
\node[cmapc] (f)  at (4.15, 0)   {knowledge suppliers as error processes};
\node[cmap]  (co) at (8.3, 0)    {commission and omission, never summed};
\node[cmap]  (b)  at (12.45, 0)  {compelled-edge bridge, per supplier};
\node[cmapr] (v)  at (0,-1.7)    {validity does not improve (S6)};
\node[cmap]  (e)  at (4.15,-1.7) {scale, 7B to 72B};
\node[cmapr] (c)  at (8.3,-1.7)  {cascade justification withdrawn};
\node[cmapg] (s)  at (12.45,-1.7) {no recitation: source named 0 of 29};
\draw[cedge] (q) -- node[above] {one of them} (f);
\draw[cedge] (f) -- node[above] {scored by} (co);
\draw[cedge] (co) -- node[above] {bias test} (b);
\draw[cedge] (f) -- node[left] {varied by} (e);
\draw[cedge] (e) -- node[above] {more claims} (v);
\draw[cedge] (e) -- node[above] {7B ties} (c);
\draw[cedge] (f.south east) -- node[above, sloped, pos=0.75] {probed by} (s.north west);
\end{tikzpicture}
\end{center}
